\clearpage
\fancyhead[L]{\small APPLICATION OF THE CALCULUS OF PROBABILITIES}
\pdfbookmark[0]{Posthumous papers. Application of the calculus of probabilities to determining the balance sheet of widows' funds}{witwen}
\SourceLeaf{124}
\begin{center}
{\Large POSTHUMOUS PAPERS.\par}\bigskip
{\large [APPLICATION OF THE CALCULUS OF PROBABILITIES\par}
\medskip{\Large TO DETERMINING THE BALANCE SHEET\par}\medskip
{\large OF WIDOWS' FUNDS ]\par}
\medskip\rule{.15\linewidth}{.3pt}\par\bigskip
{\large [I.]\par}
\pdfbookmark[1]{I. General overview of the method}{witwen-I}
\textit{[General overview of the method.]}\par
\textit{[Extract from an opinion submitted in the written vote in the University Senate.]}\par
\end{center}
The preceding [placed in brackets here] opinion of Professor D.: [The Royal University Board of Curators seems to fear that, with the large number of widows now living, the fund will sooner or later be unable to meet the pensions, which have now risen to 250 Thl. It therefore requests a report on whether there are sound grounds for such concern, and by what means any threatened danger might be averted. .\,.\,.\,] expresses the actual \textit{point at issue} so aptly that I can only assent word for word to the first half of this opinion. If doubt has been raised as to whether the fund will be able to meet its obligation on a lasting basis, this is surely the most unsuitable moment to make \textit{greater} demands on it.

\indent I can remember public opinion of this institution going back as far as 40 years. Even then it was regarded as a splendid jewel of the University, unique of its kind, precisely because of its distinctive features. \textit{Complete} freedom to join or not, and indeed to rejoin, at a comparatively small sacrifice, after having withdrawn; a very small contribution. And at that time the pension amounted to only 150 or 160 Thl. The attraction was not the size of the pension, but the \textit{generous way in which} anyone who could become a professor at Göttingen was offered secure support for a widow left behind, with the prospect of its being gradually increased under wise administration. Anyone who wished for more also joined another widows' fund. The pension has now risen to 250 Thl., and the generous manner of the institution has remained the same to this day. God grant that it may never become necessary to change \textit{anything at all} about this! Compulsory percentage deductions from salaries, intended to obtain by turning the hour hand what only the gradual progress of the minute hand can provide, would not only be far too ineffective to achieve the purpose, but would entirely destroy the noble character of the institution.

\par\SourceLeaf{125}
I am therefore of the opinion that all the changes proposed by University Councillor O. should for the moment be left entirely aside; the P. M., for which we owe him our warm thanks, has shown that there is no immediate danger. Indeed, even if the addition of more widows were to produce some deficit in the next few years in place of the latest surplus, which is still gratifying, we must not forget that the accumulated surpluses are intended in part to cover just such fluctuations arising from temporary circumstances.

But I do consider a \textit{thorough} investigation necessary, in agreement with the rescript and with the views expressed by His Magnificence. Even on the most cheerful view one might take of the society's condition, such an investigation must in any event become necessary at least later on, if only for the following reason.

If I may express my \textit{opinion} \textit{before} a thorough investigation based on rigorous calculation has taken place, I believe that the present large number of widows must be regarded as anomalous. It is true that the number of participating professors must bear a certain relation to the number of widows, and that the former number is now much greater than it used to be. But the \textit{present} high number of widows has no connection with that at all. If the number of participating professors remains as large from now on, this is \textit{a point requiring very serious consideration}, but for the distant future, not for the present; only after 20 or 30 years or more can the consequences \textit{of this} become very apparent.

On this supposition, it is probable that, perhaps within a few years, the number of widows will fall again, perhaps considerably, and that the accumulated surpluses of 8149 Thl. will therefore grow again, perhaps considerably. \textit{But whether the number of widows will fall within, for example, \textup{10} years to or below \textup{15} is much more uncertain.} Suppose now that the surpluses had grown to 15000 Thl. or more, but the number of widows stubbornly remained at 16: what is to happen then? On the one hand, one does not wish to let the surpluses grow indefinitely; on the other, the \textit{statute} prevents an increase in the pension. A thorough examination \textit{must} then be made of whether, and \textit{to what extent}, one may \textit{alter the statute} without endangering the society.

I do not wish to decline the trust with which His Magnificence and some of my colleagues honour me in wishing me to undertake such an examination, namely to draw up a balance, based on the laws of mortality and the calculus of probabilities, between the institution's assets and its obligations. I must, however, make the following preliminary remarks.

First, those gentlemen who believe that such calculations can be completed within four weeks have a very mistaken idea of how time-consuming they are. I can therefore all the less commit myself to a definite deadline, the smaller the part of my time that I shall be able to devote to them.

Second, the calculations cannot be carried out thoroughly at all without the necessary \textit{data}, of which \textit{nothing whatsoever} is available at present. I shall indicate below what the required data consist of; without them I can undertake nothing at all. Whether, how, and how soon they can be brought together I must leave entirely to the Church Deputation.

Third, the distinctive arrangements of our widows' fund contain several elements that are independent of the law of mortality and cannot be subjected to calculation. For this reason the final result must remain subject to some imperfection; I hope, however, that substitutes can be found whose use will make this imperfection insignificant.

I shall now try to indicate briefly what this work involves.

\par\SourceLeaf{126}
The obligations of the W. K. fall under three main headings.

I. Obligations towards the widows now living and, where applicable, their children,

II. Obligations towards those professors who are now members of the W. K. society.

III. Obligations towards professors who will join in future.

Ad I. The obligation towards each individual widow without minor children has exactly the value of a life annuity for her, and can therefore be calculated exactly if her present age is known (and a definite choice is made of mortality table and interest rate). That, in each individual case, a provider who took on the obligation at this price would be playing a kind of game of chance goes without saying (and I shall therefore omit this reminder below, where it would have to be made repeatedly, since it is known to anyone with some mathematical education); but a provider who entered into the same contract with a \textit{very large number} of people would, if the calculation were correct, have to expect with moral certainty only fluctuations insignificant in proportion to the whole.

Where there are children, the calculation requires a modification, but retains the same validity as in the preceding case. Naturally, the children's ages must also be known.

Finally, strictly speaking, our widows' fund, which excludes widows who remarry, would require a further modification, which is apparently\footnote{This is not the place to justify why I consider this arrangement only \textit{apparently} advantageous, but I am quite willing to give the reasons to anyone interested.} to the fund's advantage, but which of course cannot be calculated in advance and must in any case be set practically \(=0\).

Thus the total amount of I. can also be valued in money, and a part of the fund's capital assets equal to this amount must be regarded as absorbed by it.

Ad II. In a similar way, the obligation towards each individual professor who is now a married member of the society could also be valued in money if matters in our society were exactly as in those voluntary societies where the annual contribution or entrance payment is regulated according to the ages of the married couple joining. What distinguishes such a society from ours, however, is the following.

A. In the former, the contract expires if the wife dies before the husband; if it is to be renewed when he remarries, this is an entirely new contract to be regulated according to the age of the second wife. We also have unmarried members, who can choose a bride of any age, as well as widowers who may possibly remarry. None of this can now be subjected to advance calculation. I would think, however, that, if one calculated for those members who are now married on the basis of their own ages and those of their present wives, and then used the mean of these first results for the others who are not now married, matters would more or less balance out. Possibly, for some members now married, it will not be their present wives but second or third wives who will one day receive the widow's pension; on the other hand, however, a proportion of those now unmarried will undoubtedly remain so. At least I do not see what more one could do here than let the two kinds of contingency, one disadvantageous and the other advantageous to the fund, cancel one another.

B. There is also the further difference that children, perhaps not even born yet, may subsequently share in the benefits. This too cannot therefore be assessed in this way; I believe, however, that an entirely adequate substitute can be found for this imperfection, which I shall not develop further here, so as not to become too lengthy.

